\chapter{周期表不是元素通讯录}

\begin{kaichang}
拉开你家的抽屉，大概率能找到这样几样东西：一节电池、一卷铝箔、一把不锈钢勺子、一包食盐、一枚钉子。再看看窗外：树叶是绿的，天空是蓝的，你呼出的气里有二氧化碳。

现在做个小游戏。这些东西里，哪些“原料”是相同的？铝箔和易拉罐都是铝做的，这不难猜。可是：不锈钢勺子里的铁、血液里让血呈红色的铁、钉子生锈后的铁锈，是同一种“铁”吗？食盐里的氯和游泳池消毒水里的氯，是同一种东西吗？绿叶里的镁和你吃的补镁片里的镁呢？

先猜一个答案，写在纸上。这一章结束时，我们会回来对答案——而且要用一张你多半已经见过、但很可能“见过却没看懂”的表：元素周期表。
\end{kaichang}

\section{一张挂错位置的“通讯录”}

很多人第一次见到周期表，是在化学教室的墙上：一百多个小格子，每个格子里一个符号、一个数字，花花绿绿。它看起来就像一本通讯录——氢住 1 号，氦住 2 号，氧住 8 号，金住 79 号。如果真是这样，那学好化学的办法就只剩一种：背下来。

但请你凑近看一眼它的形状。这不是一张普通的方表格。它右边缺了角，底部悬空两行，像一座造型古怪的城市。为什么第一行只有 2 个格子，第二、三行各 8 个，接下来变成 18 个、32 个？为什么最下面两排——镧系和锕系——要被“流放到”表格之外？

如果周期表只是通讯录，这些形状完全没有必要。通讯录是长方形的。而周期表这副怪模样，其实是大自然亲手画的一张地图：它在告诉你，\textbf{元素的“性格”是循环出现的}。

什么叫性格循环？看看这三位：锂（\chem{Li}）、钠（\chem{Na}）、钾（\chem{K}）。它们都是银白色、软得能用刀切开的金属，都要泡在煤油里保存，遇水都会剧烈反应——而且一个比一个暴躁：锂只是嘶嘶作响地打转，钠会熔成小球四处乱窜，钾直接燃起紫色火焰。三种元素，脾气像一个模子刻出来的。在周期表上，它们恰好排在同一列。

再看最右边一列：氦、氖、氩、氪、氙、氡。这一家的性格完全相反——什么都不爱干。氦气球里的氦、霓虹灯里的氖、焊接保护气里的氩，全都不跟别的元素成键，独来独往。

同一列一个脾气，这就是周期表最大的秘密。可为什么会这样？元素又不会开会商量。答案不在表格里，在原子内部——在第 9 章认识的电子身上。

\section{电子住的“楼层”决定一切}

回忆一下第 9 章的图景：原子核外的电子不是乱飞的，它们按能量高低住在一层层“楼层”上，每层能住的人数有严格规定。第一层最多住 2 个，第二层最多 8 个，第三层最多 18 个……电子从低层往高层依次入住，就像客人从一楼开始住旅馆。

而元素的化学性格——爱不爱反应、爱跟什么反应、能交出还是抢走电子——几乎完全由\textbf{最外面那层住了几个电子}决定。内层电子被外层挡得严严实实，基本不露面。

现在，把周期表横着读，奇迹就出现了：

\begin{itemize}[itemsep=2pt]
\item 氢和氦，电子只够填第一层（2 个名额）——所以第一行只有 2 格。
\item 锂到氖，电子开始填第二层（8 个名额）——第二行恰好 8 格。
\item 钠到氩，填第三层的外圈 8 个名额——第三行也是 8 格。
\item 再往下，每层能住的电子更多，表格就随之变宽：18 格、32 格。
\end{itemize}

竖着读呢？同一列的元素，最外层电子数相同。锂、钠、钾最外层都只有孤零零的 1 个电子，所以都想“脱手”，脾气同样火爆；氦、氖、氩最外层恰好住满，所以无欲无求。表格的宽度、行数、列的家族相似性——全部来自电子楼层的容纳规则。

\begin{wuqu}
“元素反应是为了凑够 8 个电子（八隅体）。”——这是流传最广的误解之一。原子没有愿望，不会“想要”稳定。真实的说法是：当原子通过交出、抢来或共享电子，使整个体系的总能量降低时，变化就可能发生。“外层 8 电子”常常对应能量较低的排布（稀有气体恰好如此），所以它是个有用的经验口诀，但不是自然界的法律。事实上，大量稳定分子和离子并不满足八隅体——后面的章节会不断遇到例外（第 17、19 章）。把口诀当成能量规律的“速记符号”，而不是原子的“人生目标”，你就不会被骗。
\end{wuqu}

\section{门捷列夫的豪赌}

\begin{zhengju}
1869 年，俄国化学家门捷列夫在桌前摆弄几十张卡片，每张写着一种已知元素和它的性质。当时已知的元素只有 63 种，没有电子、没有原子核的概念——电子要到 1897 年才被发现。他手上唯一的线索是：按原子量从小到大排，元素的性质会\textbf{周期性地重复}。

门捷列夫做了一件当时看来近乎鲁莽的事。第一，他在表格里留了空位，宣布“这里有一种还没被发现的元素”，并且根据上下左右邻居的性格，预言了它的密度、熔点、化学性质。他给其中一个空位起名“类铝”。第二，有几处原子量顺序和性格对不上（比如碲和碘），他坚持按性格排，宁可怀疑原子量测错了。

十六年后，1875 年，法国化学家布瓦博德朗发现了新元素镓（\chem{Ga}）——位置、密度、性质，与“类铝”的预言惊人吻合。1886 年发现的锗（\chem{Ge}）又命中了“类硅”的预言。一张表居然能预言还没见过的东西，这是科学理论最硬气的时刻。顺便说，碲和碘的问题后来也水落石出：门捷列夫按性格排是对的，真正决定顺序的是原子核里的质子数（原子序数），而不是原子量——这正是第 7 章讲过的“元素身份证”。

这个故事最该记住的不是“门捷列夫真聪明”，而是方法：他相信重复出现的规律背后必有共同原因，并用空白格子把自己的信念变成可以被推翻的预言。证据来了，规律站住了，哪怕他当时说不出原因——原因（电子排布）要等半个世纪后的量子力学才补上。
\end{zhengju}

\section{表格的街区地图}

现在的周期表可以当成一座城市来逛。记住几个街区，你就不会迷路：

\begin{itemize}[itemsep=2pt]
\item \textbf{最左边两列}（第 1、2 族，碱金属和碱土金属）：最外层只有 1～2 个电子，是“急着脱手电子”的街区，性格活泼，遇水、遇酸都容易反应。
\item \textbf{中间那一大片矮房子}（过渡金属）：铁、铜、锌、钛、铬、锰都住这里。它们性格稳重又多变——能呈现好几种化合价，离子常常五颜六色（第 15 章专门逛这个街区）。
\item \textbf{右边倒数第二列}（第 17 族，卤素）：氟、氯、溴、碘，最外层 7 个电子，差一个就住满，是“见电子就抢”的街区，抢到手就变成安稳的负离子——食盐里的氯离子就是氯气抢完电子后的样子。
\item \textbf{最右边一列}（第 18 族，稀有气体）：外层住满，万事不理。
\item \textbf{底部悬空的镧系、锕系}：它们的差别藏在更深的内层电子上，外层几乎一样，所以 14 个元素性格极其相似。若把它们插回主表，表格会宽得没法印刷，于是被安排到楼下“附楼”。
\end{itemize}

还有一条看不见的斜线：从硼到砹附近，左边是金属（导电、延展、爱失电子），右边是非金属，斜线边上的硅、锗等“类金属”两边都沾点——你的手机芯片就靠硅这种“半金属半非金属”的脾气工作。

\begin{qiaoliang}
同一族往下走，为什么反应性有规律地变化？拿第 1 族来说：锂、钠、钾都是丢掉最外层那 1 个电子，凭什么钾最凶？关键在距离和遮挡。钾的那颗外层电子住在第四层，离原子核更远，中间还隔着三层内层电子像棉被一样“屏蔽”了核的吸引力（第 11 章会细讲屏蔽效应）。吸得松，就容易丢；越容易丢，遇水反应越猛。用第 27 章的能量语言说：反应放出的净能量、启动反应所需的能量，都在顺着“楼层高度”规律地变。周期表上的每一个箭头趋势，背后都是“核的吸引力”和“电子间的排斥、遮挡”两股力量的拉锯——趋势不用背箭头，理解拉锯就够了。
\end{qiaoliang}

\section{回到开场的抽屉}

现在可以兑现承诺了。不锈钢勺子、血液、铁锈里的铁：是同一种元素铁（\chem{Fe}），住在中间街区。它在勺子里是金属原子，在血液的血红蛋白里是抓住氧气的配位中心（第 15 章），在铁锈里是把电子交给了氧的氧化态——同一种元素，因为电子状态和邻居不同，扮演着完全不同的角色。

食盐里的氯和消毒水里的氯：也是同一种元素，但一个是抢到了电子、安稳度日的氯离子 \chem{Cl^-}，一个是还没抢到、脾气暴躁的氯气分子 \chem{Cl_2}。元素的“身份证”（质子数）没变，性格和危险程度却天差地别——这就是为什么“吃盐”和“吸氯气”完全是两回事。

绿叶里的镁和补镁片里的镁：同一种元素，都坐在第 2 族。在叶片里，它恰好卡在叶绿素分子的正中心，是光合作用这台机器的一颗关键螺丝钉（第 54 章）。

你看，元素从来不以“纯性格”出现在生活里，它们总是处在某种组合和电子状态中。周期表给你的不是“什么东西有毒什么无毒”的名单，而是理解一切物质行为的底牌。

\section{这座城还没逛完}

这一章我们只看了地图的轮廓。第 11 章会逐条拆解表格里那些箭头趋势——原子半径、电离能、电负性——看它们怎样从电子楼层的拉锯战里长出来；第 12 到 15 章会挨个街区拜访：暴躁的第 1 族、贪婪的第 17 族、生命骨架元素氧硫氮磷、五彩斑斓的过渡金属。而地图本身的边界也还在扩张：人工合成的新元素仍在被加到表格末尾，它们的“楼层”已经高到开始违背一些我们熟悉的规律——那正是科学还活着的证据。

\begin{tiaozhan}
为什么第二层住 8 个、第三层能住 18 个？这不是人为规定。量子力学给出四个量子数：主量子数 $n$ 决定楼层，角量子数 $l$ 决定户型（s、p、d、f），磁量子数决定朝向，自旋量子数让每个“房间”最多住两个自旋相反的电子（泡利不相容原理）。第 $n$ 层的总容量是 $2n^2$：$n=1$ 得 2，$n=2$ 得 8，$n=3$ 得 18——周期表每一行的长度，就是这些数字的直接写照。不过电子入住的顺序有个小别扭：第四层的 s 房间比第三层的 d 房间更“便宜”（能量更低），所以电子会先去 4s 再回头填 3d——这就是第 4 行突然多出 10 格过渡金属的原因，也是铬、铜等少数元素电子排布“不守规矩”的根源。细节会在第 9 章的挑战层和第 11 章继续展开。
\end{tiaozhan}

\section*{练一练}

\begin{sikaoti}
\item 不看周期表，只根据“最外层电子数决定性格”，预测：最外层有 2 个电子的镁（\chem{Mg}），和最外层有 6 个电子的氧（\chem{O}），相遇时谁更可能交出电子、谁更可能抢电子？生成物里它们各带什么电荷？
\item 氦的最外层也是“住满”的（第一层 2 个即满），氢最外层只有 1 个电子。从楼层规则解释：为什么氢气可燃，氦气却能安全充进气球？
\item 门捷列夫把碲排在碘前面，尽管碲的原子量更大。他依据的是什么？后来证明他正确的关键发现是什么？
\item 画一幅示意图：用“楼层和房间”的表示法画出钠（11 个电子）和氯（17 个电子）的电子入住情况，并用箭头标出“谁的一层快住满了、谁的一层刚开张”。图旁注明：这是帮助思考的表示法，电子并不真的住在小格子里。
\item 氟在第 17 族最顶端，氯在它楼下。参照第 1 族“楼层越高越容易丢电子”的逻辑，反方向推理：同族越往上，抢电子的能力会怎样变化？查一下周期表确认你的推理。
\end{sikaoti}
